\documentclass[11pt]{article}
\usepackage[margin=1in]{geometry}
\usepackage[T1]{fontenc}
\usepackage{lmodern}
\usepackage{microtype}
\usepackage{amsmath,amssymb,amsthm,mathtools}
\usepackage{booktabs}
\usepackage{enumitem}
\usepackage{xcolor}
\usepackage[hidelinks]{hyperref}
\usepackage{fancyhdr}

\definecolor{notegray}{RGB}{245,246,248}
\newtheorem{theorem}{Theorem}
\newtheorem{proposition}[theorem]{Proposition}
\newtheorem{lemma}[theorem]{Lemma}
\newtheorem{corollary}[theorem]{Corollary}
\newcommand{\permanent}{\operatorname{per}}
\newcommand{\R}{\mathbb{R}}
\newcommand{\N}{\mathbb{N}}
\newcommand{\one}{\mathbf{1}}

\setlength{\headheight}{14pt}
\pagestyle{fancy}
\fancyhf{}
\lhead{Dittert conjecture for $n=4$}
\rhead{Exact computer-assisted proof note}
\cfoot{\thepage}

\title{An Exact Computer-Assisted Proof of the\\Dittert Conjecture for $n=4$}
\author{Pedro Paulo Marques do Nascimento}
\date{22 July 2026\\[0.4em]
\small Licensed under \href{https://creativecommons.org/licenses/by/4.0/}{CC BY 4.0}}

\hypersetup{
  pdftitle={An Exact Computer-Assisted Proof of the Dittert Conjecture for n=4},
  pdfauthor={Pedro Paulo Marques do Nascimento},
  pdfsubject={Exact working proof and reproducible certificate},
  pdfkeywords={Dittert conjecture, permanent, computer-assisted proof, exact certificate}
}

\begin{document}
\maketitle

\begin{center}
\fcolorbox{black!25}{notegray}{\parbox{0.90\textwidth}{\small
\textbf{Review status.} This document reports a new exact computational certificate produced in an interactive research session. The certificate and three exact verifiers are included, two of which independently expand the saved factors as literal squares. The argument has not yet undergone independent peer review; it should be checked before being cited as an established result.
}}
\end{center}

\begin{abstract}
For a nonnegative $4\times4$ matrix $A$ whose entries sum to $4$, let
\[
 \phi(A)=\prod_{i=1}^{4}r_i+\prod_{j=1}^{4}c_j-\permanent(A),
\]
where $r_i$ and $c_j$ are the row and column sums. We give an exact computer-assisted proof that
\[
 \phi(A)\le 2-\frac{4!}{4^4}=\frac{61}{32},
\]
with equality only at the matrix $J_4$ whose entries are $1/4$. Published structural results reduce the boundary case to a single face with two zero entries in distinct rows and columns. On that face, an integer sum-of-squares/copositive certificate proves a strict quartic inequality. The final verification uses arbitrary-precision integer arithmetic only; the smallest residual coefficient is
\[
 \frac{72\,694\,203\,872}{2^{60}}>6.28\times10^{-8}.
\]
\end{abstract}

\section{Statement}
Let
\[
 K_4=\left\{A=(a_{ij})\in\R_{\ge0}^{4\times4}:\sum_{i,j}a_{ij}=4\right\}.
\]
Write $J_4=(1/4)_{i,j}$. The four-dimensional Dittert conjecture is the following statement.

\begin{theorem}\label{thm:main}
For every $A\in K_4$,
\[
 \phi(A)\le\frac{61}{32}.
\]
Equality holds if and only if $A=J_4$.
\end{theorem}

The proof combines two published structural facts with one exact polynomial certificate. Hwang proved that the only possible \emph{positive} $\phi$-maximizer is $J_n$ \cite{Hwang1986}. Cheon and Wanless proved that partly decomposable matrices cannot maximize $\phi$, and that a maximizer cannot have its complete zero set equal, up to row and column permutations, to one proper rectangular block \cite[Theorems~3.2 and~3.7]{CheonWanless2012}.

\section{Reduction to one two-zero face}

\begin{lemma}\label{lem:zeros}
Let $Z$ be a nonempty set of positions in a matrix. If $Z$ contains no two positions lying in distinct rows and distinct columns, then all positions of $Z$ lie in one row or all lie in one column. In particular, $Z$ is one rectangular block.
\end{lemma}

\begin{proof}
Choose $(i,j)\in Z$. Every other position in $Z$ must share row $i$ or column $j$ with $(i,j)$. If both a position $(i,j')$ with $j'\ne j$ and a position $(i',j)$ with $i'\ne i$ existed, those two positions would lie in distinct rows and distinct columns. Hence only one of the two types can occur, and all of $Z$ lies in row $i$ or in column $j$.
\end{proof}

\begin{proposition}\label{prop:reduction}
To prove Theorem~\ref{thm:main}, it is enough to prove
\[
 \phi(A)<\frac{61}{32}
 \quad\text{whenever}\quad
 A\in K_4,\qquad a_{11}=a_{22}=0.
 \tag{1}\label{eq:twozero}
\]
\end{proposition}

\begin{proof}
The compactness of $K_4$ and continuity of $\phi$ give a global maximizer $A$. Since $J_4\in K_4$, its maximum value is at least $\phi(J_4)=61/32$.

If $A$ is positive, Hwang's positive-support theorem gives $A=J_4$. Suppose instead that $A$ is on the boundary. If its zero set contained no two positions in distinct rows and columns, Lemma~\ref{lem:zeros} would put all zeroes in one row or one column. If that whole row or column were zero, the matrix would be partly decomposable; otherwise its complete zero set would be one proper rectangular block. Both alternatives contradict Cheon--Wanless. Thus it contains two zeroes in distinct rows and columns. Independent row and column permutations put those zeroes at $(1,1)$ and $(2,2)$ without changing $\phi$. Inequality \eqref{eq:twozero} would then give $\phi(A)<61/32$, contradicting maximality. Therefore the maximizer is positive and hence equals $J_4$.
\end{proof}

\section{Quartic formulation of the face}
Let $X=A/4$, so the entries of $X$ sum to $1$. On the face $x_{11}=x_{22}=0$, index the remaining fourteen entries by
\[
 V=\{(i,j):1\le i,j\le4\}\setminus\{(1,1),(2,2)\}.
\]
Write $x_v$ for $v\in V$ and $S=\sum_{v\in V}x_v$. Let $\mathcal E$ be the collection of four-element subsets $E\subseteq V$ that meet all four rows or meet all four columns. There are exactly $274$ such sets. Define
\[
 F(x)=\sum_{E\in\mathcal E}\prod_{v\in E}x_v.
 \tag{2}\label{eq:F}
\]

\begin{lemma}\label{lem:hypergraph}
On this face,
\[
 \prod_{i=1}^{4}R_i+\prod_{j=1}^{4}C_j-\permanent(X)=F(x),
\]
where $R_i$ and $C_j$ are the row and column sums of $X$. Consequently,
\[
 \phi(A)=256F(x).
 \tag{3}\label{eq:scale}
\]
\end{lemma}

\begin{proof}
Expanding the product of row sums gives every square-free quartic choosing one cell from each row. Expanding the product of column sums gives every square-free quartic choosing one cell from each column. A monomial appearing in both expansions chooses one cell from every row and every column, so it is a permanent term and is subtracted once. Therefore each monomial in the union appears exactly once, proving \eqref{eq:F}. Scaling every entry by $4$ multiplies both fourfold products and the permanent by $4^4=256$.
\end{proof}

By \eqref{eq:scale}, the required strict inequality is equivalent to
\[
 G(x):=\frac{61}{8192}S^4-F(x)>0
 \quad\text{for}\quad x\ge0,\ S=1.
 \tag{4}\label{eq:G}
\]

\section{The exact certificate}
Let $m_2(x)$ denote the vector of all $105$ degree-two monomials in the fourteen variables. The attached certificate contains:
\begin{itemize}[nosep]
 \item a group $\Gamma$ of $16$ permutations of the variables preserving the two-zero face and $F$;
 \item one representative from each of the $16$ orbits of unordered variable pairs under $\Gamma$;
 \item an integer lower-triangular matrix $L_0\in\mathbb Z^{105\times105}$;
 \item integer lower-triangular matrices $L_p\in\mathbb Z^{14\times14}$ for the $16$ pair representatives.
\end{itemize}
All actual Cholesky factors are $2^{-30}L$. The integer-only verifier expands the following identity coefficient by coefficient:
\begin{align}
 2^{60}G(x)
 &=\sum_{\gamma\in\Gamma}
      \left\|L_0^{\mathsf T}m_2(\gamma x)\right\|_2^2 \notag\\
 &\quad+\sum_{p=(u,v)}\sum_{\gamma\in\Gamma}
      x_{\gamma u}x_{\gamma v}
      \left\|L_p^{\mathsf T}(\gamma x)\right\|_2^2
   +\sum_{|\alpha|=4}r_\alpha x^\alpha.
 \tag{5}\label{eq:certificate}
\end{align}
Here every $r_\alpha$ is an integer and
\[
 \min_{|\alpha|=4}r_\alpha=72\,694\,203\,872>0.
 \tag{6}\label{eq:minres}
\]
There are
\[
 \binom{14+4-1}{4}=2380
\]
quartic monomials, and the verifier checks every one.

\begin{proposition}\label{prop:face}
The strict inequality \eqref{eq:G} holds on the nonnegative simplex $S=1$.
\end{proposition}

\begin{proof}
Identity \eqref{eq:certificate} is an exact polynomial identity over the integers. Its first sum is a sum of squares. Every summand in the second sum is nonnegative when $x\ge0$, because it is a product of two nonnegative coordinates and a square norm. The final polynomial has strictly positive coefficients by \eqref{eq:minres}. Since $S=1$, at least one coordinate is positive, and the corresponding fourth-power monomial makes the final polynomial strictly positive. Hence $G(x)>0$.
\end{proof}

Proposition~\ref{prop:face}, Lemma~\ref{lem:hypergraph}, and Proposition~\ref{prop:reduction} prove Theorem~\ref{thm:main}.

\section{Exact verification and reproducibility}
The numerical semidefinite search was used only to \emph{discover} the factors. It is not part of the proof. The factors were rounded to integers, after which all proof decisions were made with arbitrary-precision integer arithmetic.

Three exact verification programs are supplied:
\begin{enumerate}[nosep]
 \item The primary verifier forms each exact Gram matrix $LL^{\mathsf T}$, expands the invariant polynomial, validates the symmetry group and all $274$ hyperedges, and checks all $2380$ residual coefficients.
 \item The NumPy literal-square audit never forms a Gram matrix. It expands every saved factor literally and constructs the target directly from $S^4$, the row and column products, and all $24$ permanent permutations.
 \item The third audit repeats the calculation without NumPy and parses the certificate archive itself.
\end{enumerate}
All three return the same minimum coefficient in \eqref{eq:minres} and terminate with \texttt{CERTIFIED}. The certificate SHA-256 digest is
\begin{center}
\texttt{d76533bd1c5566ea8d96aa3b58a0b6a8bf3310eb438776ebd99a6bd88d5a11f6}.
\end{center}

Run the checks from the bundle directory, without Python's \texttt{-O} option:
\begin{verbatim}
python verify_primary.py dittert_n4_exact_certificate.npz
python verify_literal_square_numpy.py dittert_n4_exact_certificate.npz
python verify_literal_square_stdlib.py dittert_n4_exact_certificate.npz
\end{verbatim}

\section{Scope and literature context}
This certificate settles only the case $n=4$. Separate July 2026 working material distributed with this note treats $n=14,15,16$; Pang's June 2026 preprint treats all $n\ge17$ \cite{Pang2026}, and Kafidov independently treats $n=16$ \cite{Kafidov2026}. No claim about dimensions $5$ through $13$ is made here. The published 2024 paper of Udayan and Somasundaram proves that a maximizing matrix in $K_4$ is fully indecomposable, but that stronger structural result is not needed here \cite{UdayanSomasundaram2024}. The present reduction uses only Hwang's positive-support theorem and Cheon--Wanless's two boundary exclusions.

\begin{thebibliography}{9}
\bibitem{Hwang1986}
S.-G. Hwang,
\emph{A note on a conjecture on permanents},
Linear Algebra and its Applications 76 (1986), 31--44.

\bibitem{CheonWanless2012}
G.-S. Cheon and I. M. Wanless,
\emph{Some results towards the Dittert conjecture on permanents},
Linear Algebra and its Applications 436 (2012), 791--801.

\bibitem{UdayanSomasundaram2024}
D. K. Udayan and K. Somasundaram,
\emph{Lih Wang's and Dittert's conjectures on permanents},
Special Matrices 12 (2024), article 20240006.

\bibitem{Pang2026}
Z. Pang,
\emph{Proof of Dittert's conjecture for dimensions $n\ge17$},
arXiv:2606.01531 (2026), preprint.

\bibitem{Kafidov2026}
B. Kafidov,
\emph{Dittert's conjecture in dimension 16 via a joint-deficit scaling lemma},
arXiv:2607.19439 (2026), preprint.
\end{thebibliography}

\end{document}
